\section{Method}
\label{sec:method}

% 3.1 The FK training objective. Give the estimator, state whether the target is
%     detached, and prove or state its unbiasedness. Be explicit that
%     E[(u - Yhat)^2] = (u - E Y)^2 + Var(Yhat)/M, so the loss has a noise floor at
%     Var/M -- a reader who sees the loss plateau above zero must already know why.
% 3.2 Path features. Endpoint vs. full signature vs. log-signature. Give the feature
%     dimension explicitly (sum_k d^k) and confront its growth in d.
% 3.3 Variance reduction. Control variate with the network's own value as control;
%     antithetic pairing; importance sampling with ESS monitoring.
% 3.4 Algorithm box: one pseudocode listing of a single FK-PINN training step.
% 3.5 Cost model. Per-step cost of \pinn{} vs \fkpinn{} as a function of d, M and K.
%     This is what makes the equal-wall-clock comparison in Sec.~\ref{sec:experiments}
%     interpretable rather than merely empirical.

Placeholder.
